\section{The central band}\label{sec:central}

We retain the notation of Section~\ref{sec:continuation}.  In the band
containing the signed coordinates, we sharpen the basic energy comparison
to preserve their exact edge cost.  The resulting bound will force any
function satisfying the crossing conditions to be a signed coordinate.

\begin{proposition}\label{prop:central}
Let $f$ be a nonconstant Boolean function, let $0<\rho<1$, and suppose that
$s\ge .84$ and $0\le m\le r^2/2$.  If
\begin{equation}\label{cen:crossing}
 b=h-1+s,\qquad \mathcal E(d)\le K=\frac rs,
\end{equation}
then $f$ is a signed coordinate.
\end{proposition}

Throughout the proof all terminating decimals denote exact rational
numbers.  Under \eqref{cen:crossing}, set
\begin{equation}\label{cen:calibration}
 y=\sqrt b,\qquad x_0=\sqrt{1-b^2},\qquad
 \alpha=\frac{F(x_0)}{x_0},\qquad M=F(x_0)^2.
\end{equation}
The symbol $\alpha$ is local to this section.  Write $R_0=.2944$.  Since
$0<r\le R_0$ and $m\le R_0^2/2$, the exact comparison
$(.999)^2<1-R_0^4/4$ gives
\begin{equation}\label{cen:domain}
 \begin{gathered}
 h>.999,\qquad .839<b\le s<1,\qquad y>.91595>.915,\\
 0\le s-b=\frac{m^2}{1+h}<.001,\qquad M\ge x_0^2\ge r>0.
 \end{gathered}
\end{equation}
Here $F(x)\ge x$ for $x\ge0$, and $(.91595)^2<.839$.
The symmetric pair of values $\pm x_0$ has $H(x_0)=b$, so its
$H$-average equals that of $d$ and its transformed squared amplitude is $M$.
Its edge energy is $x_0^2/b$.  When $m=0$, the crossing identity gives
$b=s$ and $x_0=\rho$, recovering a noised coordinate.

\subsection{The calibrated edge cost}

For $0\le v<F(1)^2$, define
\begin{equation}\label{cen:phi}
 \Phi(v)=\frac{x^2}{H(x)}-v,
 \qquad x=F^{-1}(\sqrt v),\qquad
 j(v)=\begin{cases}\Phi(v)/v,&v>0,\\0,&v=0.\end{cases}
\end{equation}
The endpoint value $F(1)$ is finite, since the integrand defining $F$ has
an integrable singularity of order $3/4$ at $1$.

\Needspace{6\baselineskip}
\begin{lemma}[Calibrated edge comparison]\label{cen:edge}
For every coordinate edge,
\begin{equation}\label{cen:edge-bound}
 D_i d\,D_i\!\left(\frac d{H(d)}\right)
 \ge (D_i g)^2+\Phi\bigl((D_i g)^2\bigr).
\end{equation}
The functions $j$ and $\Phi$ are nonnegative, increasing, and convex.
Furthermore, $\Phi'(v)/v$ is increasing for $v>0$.
\end{lemma}

\begin{proof}
Put $G=F(x)$ and $\ell(G)=\log F'(x)$.  Direct differentiation gives
\begin{equation}\label{cen:ell}
 \ell'(G)=\frac{3x}{2\sqrt{H(x)}},\qquad
 \ell''(G)=\frac34\left(H(x)+\frac1{H(x)}\right),\qquad
 \ell'''(G)=\frac{3x^3}{4H(x)^{3/2}}.
\end{equation}
Thus $\ell$ is even and convex, and $\ell''$ is increasing on the
nonnegative half of its domain.

Consider an edge with distinct $g$-values $G_-<G_+$, and let
$I=[G_-,G_+]$.  Changing variables from $x$ to $G$ in the two differences
shows that the ratio of its energy to $(D_i g)^2$ is
\[
 \left(\frac1{|I|}\int_I e^{\ell(G)}\,dG\right)
 \left(\frac1{|I|}\int_I e^{-\ell(G)}\,dG\right)
 =\E\cosh\bigl(\ell(U)-\ell(V)\bigr),
\]
where $U,V$ are independent uniform points of $I$.  For a fixed separation
$\delta=|U-V|$, the midpoint is uniform on the interval of length
$|I|-\delta$ with the same center as $I$.  The function
\[
 A_\delta(c)=\ell(c+\delta/2)-\ell(c-\delta/2)
\]
is odd, and is increasing because
$A_\delta'(c)=\ell'(c+\delta/2)-\ell'(c-\delta/2)\ge0$.
Consequently $\cosh A_\delta(c)$ is even and increasing in $|c|$.
The integral of any such function over an interval of fixed length is
minimized by centering that interval at zero: for positive center $c$,
the derivative of the integral is the value at the right endpoint minus
the value at the left endpoint, which is nonnegative.  The domain of $\ell$ is the symmetric interval $(-F(1),F(1))$.
Moving the center of $I$ toward zero while keeping its length fixed
keeps $I$, and hence every conditional midpoint interval, in this domain.
Applying the preceding fixed-length integral comparison conditionally
on $\delta$ proves that the displayed product of averages is minimized
when $I$ is centered at zero.

If that centered interval has endpoints $\pm a$, its corresponding
$d$-values are $\pm x$, where $F(x)=a$.  Its energy is exactly
$x^2/H(x)=a^2+\Phi(a^2)$.  Since $a^2=(D_i g)^2$, this proves
\eqref{cen:edge-bound}.  A zero difference gives the same statement
directly.

The symmetric-interval computation also gives, for $v>0$,
\begin{equation}\label{cen:j-average}
 j(v)=\frac14\int_{-1}^1\int_{-1}^1
 \left[\cosh\bigl(\ell(\sqrt v\,a)-\ell(\sqrt v\,c)\bigr)-1\right]
 \,da\,dc.
\end{equation}
To prove convexity of the integrand in $v$, use evenness of $\ell$ and
$\cosh$ to replace $(a,c)$ by numbers $A\ge C\ge0$.
For $B(v)=\ell(A\sqrt v)-\ell(C\sqrt v)$, differentiation gives
\[
 B'(v)=\frac{z_A\ell'(z_A)-z_C\ell'(z_C)}{2v},\qquad
 B''(v)=\frac{Q(z_A)-Q(z_C)}{4v^2},
 \quad z_A=A\sqrt v,\ z_C=C\sqrt v,
\]
where $Q(z)=z(z\ell''(z)-\ell'(z))$.
The function $z\ell'(z)$ is increasing on $z\ge0$, and
\[
 Q'(z)=z\ell''(z)-\ell'(z)+z^2\ell'''(z)\ge0.
\]
Indeed, $\ell'(0)=0$ and increasing $\ell''$ imply
$\ell'(z)\le z\ell''(z)$, and \eqref{cen:ell} gives $\ell'''(z)\ge0$.
Thus $B$ is nonnegative, increasing, and convex.  Its composition with
$\cosh-1$ has the same properties.  Averaging proves these properties
for $j$, including its continuous extension $j(0)=0$.
Now $\Phi(v)=vj(v)$ gives
$\Phi'=j+vj'\ge0$ and $\Phi''=2j'+vj''\ge0$.
Finally, convexity and $j(0)=0$ imply that $j(v)/v$ is increasing, while
$j'$ is increasing.  Therefore
\[
 \frac{\Phi'(v)}v=\frac{j(v)}v+j'(v)
\]
is increasing as claimed.
\end{proof}

Define the slope and energy excess at the calibration by
\begin{equation}\label{cen:pdelta}
 p_*:=\Phi'(M)=\frac{1+b^2}{2\alpha b^{3/2}}-1,\qquad
 \Delta:=M+\Phi(M)-K=(s-b)\left(1+\frac1{bs}\right).
\end{equation}
The first identity follows by differentiating $x^2/H(x)$ with respect
to $F(x)^2$; the second uses
$M+\Phi(M)=(1-b^2)/b$ and $K=(1-s^2)/s$.
In particular, $p_*,\Delta\ge0$.

\begin{lemma}[Calibration bounds]\label{cen:constants}
At a crossing in the central band,
\begin{equation}\label{cen:four-constants}
 \frac{\Phi(M)}M<.0063,\qquad p_*<.019,\qquad
 \frac{p_*}M<.0532,\qquad \frac\Delta M<.008.
\end{equation}
\end{lemma}

\begin{proof}
As $b$ decreases, $M$ increases.  Lemma~\ref{cen:edge} therefore reduces
the first three upper bounds to $b_0=.839$.  Set
$\alpha_0=F(\sqrt{1-b_0^2})/\sqrt{1-b_0^2}$.
At this endpoint the positive binomial series, integrated term by term,
gives
\[
 \alpha_0\ge\sum_{j=0}^{5}
 \frac{(3/4)_j(1-b_0^2)^j}{j!(2j+1)}>a_0:=1.08834,
 \qquad \sqrt{b_0}>y_0:=.91595,
\]
where $(3/4)_0=1$ and $(3/4)_j=(3/4)(7/4)\cdots(3/4+j-1)$.
Using $\Phi(M)/M=1/(b\alpha^2)-1$ and \eqref{cen:pdelta}, exact rational
comparisons give
\begin{align*}
 \frac1{b_0a_0^2}-1&<.006257<.0063,\\
 P_0:=\frac{1+b_0^2}{2a_0b_0y_0}-1&<.018641<.019,\\
 \frac{P_0}{(1-b_0^2)a_0^2}&<.053154<.0532.
\end{align*}
For the last bound, \eqref{cen:domain} and $m^2\le r^4/4$ give
\[
 \frac\Delta M\le
 \frac{r^3}{4(1+h)}\left(1+\frac1{bs}\right)
 <\frac{R_0^3}{4(1.999)}
       \left(1+\frac1{(.839)(.84)}\right)<.007720<.008.
\]
The numerical comparisons above are exact rational inequalities.
The four bounds in \eqref{cen:four-constants} are also verified by the
central-band certificate described in Appendix~\ref{app:arithmetic}.
\end{proof}

\subsection{Reduction to variance and first-level estimates}

Let $a_i$ be
the degree-one Fourier coefficients of $g$, and put
\[
 q=\sum_i a_i^2,\qquad l=\max_i a_i^2,\qquad H_0=\Var g-q.
\]
Thus $H_0$ is the Fourier weight above degree one.  Counting degrees gives
$\sum_i\E(D_i g)^2\ge\Var g+H_0$.
All arguments of $\Phi$ used below belong to its open domain:
since the cube is finite, $G_{\max}:=\max|g|<F(1)$; every half-difference
$|D_i g|$ and coefficient $|a_i|$ is at most $G_{\max}$, while
$M=F(x_0)^2<F(1)^2$ because $b>0$.
In a coordinate attaining $l$, Jensen's inequality and monotonicity of
$\Phi$ give
$\E\Phi((D_i g)^2)\ge\Phi(\E(D_i g)^2)\ge\Phi(l)$.
Lemma~\ref{cen:edge} and the tangent inequality for $\Phi$ at $M$ yield
\begin{align}
 \mathcal E(d)&\ge\Var g+H_0+\Phi(l)\label{cen:spectral-energy}\\
 &\ge K+(1+p_*)(\Var g-M)+(1-p_*)H_0
             -p_*(q-l)+\Delta.\label{cen:tangent-energy}
\end{align}

Here $q-l$ is the degree-one weight outside a largest coordinate.
Put
\begin{equation}\label{cen:V0}
 V_0=\E(g^2-M)^2,\qquad C_0=.034.
\end{equation}
We will prove
\begin{equation}\label{cen:variance-gain}
 \Var g-M\ge C_0V_0-\frac\Delta{1+p_*},
\end{equation}
with strict inequality unless $m=0$ and $g^2=M$ everywhere.
Substituting this estimate into \eqref{cen:tangent-energy} cancels
$\Delta$.  We will then control the remaining negative term by proving
\[
 p_*(q-l)\le(1+p_*)C_0V_0+(1-p_*)H_0.
\]
These two bounds give $\mathcal E(d)>K$ whenever the variance inequality
is strict.  We first prove that inequality, then establish the
first-level comparison and treat the remaining equality case.

\subsection{A variance gain around the calibration}

Define, on $0\le t<1$,
\[
 k(t)=\frac{F(\sqrt{1-t^4})}{\sqrt{1-t^4}},\qquad k(1)=1.
\]
Writing $u=\sqrt{H(d)}$, one has $g=k(u)d$, $k(y)=\alpha$, and
\begin{equation}\label{cen:k-ode}
 k'(t)=-\frac{2(1-t^3k(t))}{1-t^4},\qquad
 g^2-M=-4\int_y^u k(t)\,dt.
\end{equation}
Both identities follow by differentiation; the second also holds at
$u=1$ by continuity.

\begin{lemma}[Bounds for the comparison function]\label{cen:k-bounds}
For $0\le t\le1$,
\begin{equation}\label{cen:k-estimates}
 2-t+\tfrac12(1-t)^2\le k(t)\le2-t+\tfrac34(1-t)^2,
 \qquad
 1\le-k'(t)\le\min\{2,1+\tfrac32(1-t)\}.
\end{equation}
At $t=1$, the derivative is understood as its left limit.
\end{lemma}

\begin{proof}
For $P_c(t)=2-t+c(1-t)^2$, define the differential residual
\[
 E_c(t)=-(1-t^4)P_c'(t)-2(1-t^3P_c(t)).
\]
Expansion yields
\begin{align*}
 \frac{E_{1/2}(t)}{(1-t)^2}&=t(2t^2-2t-1)\le0,\\
 \frac{E_{3/4}(t)}{(1-t)^2}
 &=\tfrac12(2t+1)(3t^2-3t+1)>0.
\end{align*}
The first sign uses $2t^2-2t-1\le-1$; for the second,
$3t^2-3t+1=3(t-1/2)^2+1/4$.
To justify backward comparison despite the singular endpoint, fix
$t<v<1$.  The error $e=k-P_c$ satisfies
$e'-2t^3e/(1-t^4)=E_c/(1-t^4)$, whence
\[
 e(t)=\sqrt{\frac{1-v^4}{1-t^4}}\,e(v)
 -\int_t^v\sqrt{\frac{1-z^4}{1-t^4}}
                   \frac{E_c(z)}{1-z^4}\,dz.
\]
Both $k(v)$ and $P_c(v)$ tend to $1$.  The first term therefore vanishes
as $v\uparrow1$, and the integral converges.  Its sign proves the two
polynomial comparisons in \eqref{cen:k-estimates}.

Substitute these comparisons into \eqref{cen:k-ode}.  The sufficient
numerators for the three derivative inequalities are, respectively,
\begin{align*}
 2(1-t^3P_{3/4})-(1-t^4)
   &=\tfrac12(1-t)^2(-3t^3+6t^2+4t+2),\\
 2(1-t^4)-2(1-t^3P_{1/2})
   &=t^3(t-5)(t-1),\\
 \tfrac{5-3t}{2}(1-t^4)-2(1-t^3P_{1/2})
   &=\tfrac12(1-t)^2(5t^3-3t^2-t+1).
\end{align*}
The first two are nonnegative on $[0,1]$.  For the last cubic, on
$[0,.6]$ use
$5t^3-3t^2+.16=(5t+1)(5t-2)^2/25\ge0$, leaving at least $.24$.
On $[.6,1]$ write the cubic as $t^2(5t-3)+(1-t)\ge0$.
This proves the derivative bounds for $t<1$.  The Taylor expansion
$F(x)/x=1+x^2/4+O(x^4)$ at $x=0$ shows that $k'(1)=-1$.
\end{proof}

To estimate the variance of $g$, separate its mean from the known mean
$\E d=m$ by subtracting $\alpha d$.  Define
\begin{equation}\label{cen:RD}
 R=g-\alpha d,\qquad
 D(u)=g^2-M-R^2+\frac{2\alpha}{y}(u^2-y^2).
\end{equation}
The right side defining $D$ depends only on $u$, because
$R^2=(1-u^4)(k(u)-\alpha)^2$.  The added multiple of $u^2-y^2$
has zero expectation.  Since $\E u^2=b=y^2$, $\E g=\E R+\alpha m$, and
$(\E R)^2\le\E R^2$, one obtains
\begin{equation}\label{cen:variance-rearrange}
 \Var g-M\ge\E\bigl[D-2\alpha m|R|\bigr]-\alpha^2m^2.
\end{equation}

\begin{lemma}[Pointwise remainder bound]\label{cen:pointwise}
Under the central-band crossing hypotheses of Proposition~\ref{prop:central},
for every $0\le u\le1$,
\begin{equation}\label{cen:pointwise-bound}
 D(u)-2\alpha m|R|+.88m^2
 \ge C_0(g^2-M)^2+.20(u^2-y^2)^2.
\end{equation}
The inequality is strict whenever $u\ne y$.
\end{lemma}

\begin{proof}
The integral identity in \eqref{cen:k-ode} gives exactly
\begin{equation}\label{cen:D-integral}
 D(u)=\frac{2\alpha}{y}(u-y)^2
       +4\int_u^y(k(t)-\alpha)\,dt
       -(1-u^4)(k(u)-\alpha)^2.
\end{equation}
Let $L=|u-y|$ and $\lambda=|k(u)-k(y)|$.  Monotonicity of $k$ makes
the integral nonnegative.  A function with slope of magnitude at most
$2$ needs length at least $\lambda/2$ to rise from $0$ to $\lambda$;
the triangle of that base and height is a lower bound for its integral.
Thus
\[
 4\int_u^y(k(t)-\alpha)\,dt\ge\lambda^2.
\]
The lower derivative bound in Lemma~\ref{cen:k-bounds} gives
$\lambda\ge L$, so \eqref{cen:D-integral} implies
\begin{equation}\label{cen:D-lower}
 D(u)\ge\left(\frac{2\alpha}{y}+u^4\right)L^2.
\end{equation}

For explicit estimates, put $a=1-u$, $c=1-y$, and
\begin{align*}
 A_-&=1+c+\tfrac12c^2,& A_+&=1+c+\tfrac34c^2,\\
 B&=1+\tfrac12(a+c)+\tfrac14(a^2+ac+c^2),&
 J&=1+\tfrac34(a+c).
\end{align*}
Lemma~\ref{cen:k-bounds}, and its averages over the interval joining
$u$ and $y$, give
\begin{equation}\label{cen:BJ}
 A_-\le\alpha\le A_+,\qquad
 |g^2-M|\le4BL,\qquad
 |R|\le J L\sqrt{1-u^4}.
\end{equation}
In fact $B$ is the average of the upper quadratic bound for $k$, and
$J$ is the average of the upper linear bound for $-k'$.
Define
\begin{equation}\label{cen:S}
 S(u,y)=2A_-+y\bigl(u^4-16C_0B^2-.20(u+y)^2\bigr).
\end{equation}
By \eqref{cen:D-lower}--\eqref{cen:BJ}, the left side minus the right
side of \eqref{cen:pointwise-bound} is at least
\begin{equation}\label{cen:pointwise-reduction}
 \frac{L^2}{y}S
 -2A_+JmL\sqrt{1-u^4}+.88m^2.
\end{equation}
If $u\le .8$, then $L=y-u>0$.  Since $b\le s$ and $2m\le r^2$,
\[
 2m\le r^2\le(1-b^2)^2=(1-y^4)^2.
\]
Dropping $.88m^2$ and bounding $\sqrt{1-u^4}\le1$ makes
\eqref{cen:pointwise-reduction} at least $LP_1/y$, where
\begin{equation}\label{cen:P1}
 P_1(u,y)=(y-u)S-A_+J(1-y^4)^2y.
\end{equation}
If $u\ge .8$ and $L>0$, the elementary inequality
\[
 2A_+JmL\sqrt{1-u^4}
 \le .88m^2+\frac{(A_+J)^2L^2(1-u^4)}{.88}
\]
makes the same expression at least $L^2P_2/(.88y)$, where
\begin{equation}\label{cen:P2}
 P_2(u,y)=.88S-(A_+J)^2(1-u^4)y.
\end{equation}
It therefore suffices, when $L>0$, to prove the polynomial signs
\begin{equation}\label{cen:polynomial-signs}
 \begin{aligned}
 P_1(u,y)&>0 &&(0\le u\le .8,\ .915\le y\le1),\\
 P_2(u,y)&>0 &&(.8\le u\le1,\ .915\le y\le1).
 \end{aligned}
\end{equation}
Here is a finite rational certificate for these signs, using the Bernstein
coefficient formula of Appendix~\ref{app:arithmetic}. Map each stated
rectangle affinely to $[0,1]^2$, with the $u$ coordinate first.  If the
resulting power coefficients are $c_{ab}$, its Bernstein coefficient of
index $(i,j)$ and order $(N_1,N_2)$ is
\begin{equation}\label{cen:bernstein-formula}
 \beta_{ij}=\sum_{a\le i,\,b\le j}c_{ab}
       \frac{\binom ia}{\binom{N_1}a}
       \frac{\binom jb}{\binom{N_2}b}.
\end{equation}
The orders and row minima are as follows; the final column records
$\lfloor1000\min_j\beta_{ij}\rfloor$ for successive $i$.
\begin{center}
\begin{tabular}{@{}ccl@{}}
\toprule
Polynomial & Order $(N_1,N_2)$ & Scaled row minima \\
\midrule
$P_1$ & $(5,12)$ & $134,353,365,278,214,46$ \\
$P_2$ & $(6,7)$ & $179,364,559,765,983,1214,1457$ \\
\bottomrule
\end{tabular}
\end{center}
Thus every coefficient is strictly positive.  Bernstein basis functions
are nonnegative and sum to one on the unit square, proving
\eqref{cen:polynomial-signs} on the closed rectangles.  The formula and
the explicit polynomials specify all arithmetic in the table; the
executable certificate is \texttt{verification/central\_band.py}.
Appendix~\ref{app:arithmetic} gives the general certificate convention.

The two lower bounds above are therefore strictly positive whenever
$L>0$.
Finally, when $u=y$, all terms in \eqref{cen:pointwise-bound} vanish
except $.88m^2$, which is nonnegative.  This also proves the asserted
strictness.
\end{proof}

\begin{lemma}[Strict variance gain]\label{cen:strict-variance}
Inequality \eqref{cen:variance-gain} holds.  If $m>0$, the stronger bound
\begin{equation}\label{cen:positive-margin}
 \Var g-M>C_0V_0-\frac\Delta{1+p_*}+.016m^2
\end{equation}
holds.  If $m=0$, \eqref{cen:variance-gain} is strict unless $g^2=M$
at every point of the cube.
\end{lemma}

\begin{proof}
Write $T=\E(u^2-y^2)^2=\E(H(d)-b)^2$.
Since $\E d=m$ and $H(d)^2=1-d^2$,
$T=h^2-b^2-\Var d$.  Lemma~\ref{lem:fourier-bounds}, using $w\le1-m$,
therefore gives
\[
 T\ge(1-r^2)h^2-b^2-rs^2(1-m).
\]
The crossing identity implies $sh-b=(1-s)(1-h)\ge0$.
Substituting $b^2\le s^2h^2$ proves
\begin{equation}\label{cen:T-lower}
 T\ge rs^2(m-m^2).
\end{equation}
For $m>0$ this yields
\begin{equation}\label{cen:T-ratio}
 \frac T{m^2}\ge s^2\left(\frac2r-r\right)
 \ge(1-R_0)\left(\frac2{R_0}-R_0\right)>4.5.
\end{equation}
Both positive factors in the middle expression decrease with $r$ on
$0<r\le R_0$, and the final inequality is rational.

We also need to compare the calibration loss with $\alpha^2m^2$.
The upper bound for $k(y)$ gives $\alpha<1.1$ and
\[
 \alpha y\le(1-c)(1+c+\tfrac34c^2)
       =1-\tfrac14c^2-\tfrac34c^3\le1,
 \qquad c=1-y\in[0,.085].
\]
Using \eqref{cen:pdelta} and $s-b=m^2/(1+h)$, for $m>0$ we obtain
\begin{align}
 \frac\Delta{(1+p_*)m^2}
 &=\frac\alpha y\,\frac2{1+h}
       \left(1-\frac{s-b}{s(1+b^2)}\right)\label{cen:delta-cancel}\\
 &>\alpha^2-\frac{(1.1)^2(.001)}{.84}
 >\alpha^2-.004.\notag
\end{align}
In the first factor of \eqref{cen:delta-cancel}, $\alpha/y\ge\alpha^2$;
also $2/(1+h)\ge1$.  The subtracted relative loss is less than
$.001/.84$, by \eqref{cen:domain}, which proves the next line.

Averaging \eqref{cen:pointwise-bound} and applying
\eqref{cen:variance-rearrange} now gives
\[
 \Var g-M\ge C_0V_0+.20T-(\alpha^2+.88)m^2.
\]
For $m>0$, insert \eqref{cen:T-ratio} and \eqref{cen:delta-cancel}.
The remaining strict margin is
$.20(4.5)-.88-.004=.016$, proving \eqref{cen:positive-margin}.
For $m=0$ one has $b=s$ and $\Delta=0$; the same averaged inequality
proves \eqref{cen:variance-gain}.  If any cube point has $u\ne y$,
Lemma~\ref{cen:pointwise} makes the averaged inequality strict.
Since $k>0$ in \eqref{cen:k-ode}, the condition $u=y$ is equivalent to
$g^2=M$, completing the strictness statement.
\end{proof}

\subsection{Concentration of the first Fourier level}

\begin{proof}[Proof of Proposition~\ref{prop:central}]
Recall that $q$ is the total degree-one Fourier weight of $g$, $l$ is
its largest coordinate weight, and $H_0=\Var g-q$ is its weight above
degree one.

The resources available to control $q-l$ follow without losing the
$\Delta$ term.  Namely, \eqref{cen:crossing},
\eqref{cen:spectral-energy}, and \eqref{cen:variance-gain} imply
\begin{equation}\label{cen:resources}
 H_0+C_0V_0\le\Phi(M)-\frac{p_*\Delta}{1+p_*}\le\Phi(M),
 \qquad
 q\ge M-\Phi(M)-\Delta>.985M>0.
\end{equation}
Indeed, substituting the variance gain into the first energy bound gives
\[
 H_0+C_0V_0\le K-M+\frac{\Delta}{1+p_*}-\Phi(l)
 =\Phi(M)-\frac{p_*\Delta}{1+p_*}-\Phi(l).
\]
Since $\Phi(l)\ge0$, this proves the first resource bound. Combining it
with the variance gain once more yields
\[
 q\ge M-\Phi(M)+2C_0V_0+\frac{p_*-1}{1+p_*}\Delta
 \ge M-\Phi(M)-\Delta.
\]
Thus the second inequality retains the full calibration loss;
Lemma~\ref{cen:constants} gives
$1-.0063-.008=.9857>.985$.

Let
\[
 L_1(x)=\sum_i a_i x_i,\qquad P=L_1^2-q,\qquad
 W=\E P^2=2\left(q^2-\sum_i a_i^4\right).
\]
Since $\sum_i a_i^4\le lq$,
\begin{equation}\label{cen:spread}
 q-l\le\frac W{2q}.
\end{equation}
Write $g=\mu+L_1+Z$, with $\mu=\E g$ and $Z$ containing only Fourier
degrees at least two; thus $\E Z^2=H_0$.
The degree-three part of $L_1P$ has coefficient $6a_i a_j a_k$ on
$x_ix_jx_k$, so its squared norm is
\[
 36\sum_{i<j<k}a_i^2a_j^2a_k^2\le3qW.
\]
To see the last inequality, write $W=4\sum_{i<j}a_i^2a_j^2$:
in $qW$ each distinct triple occurs three times, with total coefficient
$12$, and all remaining terms are nonnegative.
Only that degree-three part of $L_1P$ pairs with $Z$.
Also $P\ge-q$, $\E P=\E L_1P=0$, and $\E L_1^2P=W$.
Expanding $g^2P$ and applying Cauchy--Schwarz therefore gives
\begin{equation}\label{cen:quadratic}
 \sqrt{V_0W}\ge\E[g^2P]
 \ge W-(2\sqrt{3q}+2|\mu|)\sqrt{WH_0}-qH_0.
\end{equation}
The first inequality uses $\E[g^2P]=\E[(g^2-M)P]$.
For the last inequality the three error terms are
$2\E(L_1PZ)$, $2\mu\E(PZ)$, and $\E(Z^2P)$, bounded below by
$-2\sqrt{3qWH_0}$, $-2|\mu|\sqrt{WH_0}$, and $-qH_0$ respectively.

The mean term is small.  Lemma~\ref{cen:k-bounds} and
\eqref{cen:k-ode} give
\[
 |g-\alpha d|\le2|u-y|\le\tfrac12|g^2-M|.
\]
Together with \eqref{cen:resources} this proves
\begin{equation}\label{cen:mean-bound}
 \frac{|\mu|}{\sqrt q}
 \le\frac{m/x_0+\tfrac12\sqrt{.0063/C_0}}{\sqrt{.985}}
 \le\frac{R_0(.543)/2+.431/2}{.9924}<.298<.31.
\end{equation}
Here $m/x_0\le r^{3/2}/2$, and the rational comparisons
$(.543)^2>R_0$, $(.431)^2>.0063/.034$, and
$(.9924)^2<.985$ justify all square-root bounds.

Set $A=\sqrt{V_0}$, $B_0=\sqrt{qH_0}$, and $z=\sqrt W$.
Equations~\eqref{cen:quadratic} and \eqref{cen:mean-bound} imply
\[
 z^2\le Az+(2\sqrt3+.62)B_0z+B_0^2.
\]
Writing $c_0=2\sqrt3+.62$, the positive root of this quadratic is at
most $A+4.32B_0$.  Indeed, substituting that value in the polynomial
$z^2-(A+c_0B_0)z-B_0^2$ gives
\[
 (4.32-c_0)AB_0+(4.32^2-4.32c_0-1)B_0^2\ge0,
\]
and the substituted value is beyond its vertex.  Both coefficients
are positive: $\sqrt3<1.73206$ gives
$4.32>c_0$ and $4.32^2>1+4.32c_0$ by rational comparison.
The degenerate cases $B_0=0$ or $W=0$ obey the same bound.  Thus
\begin{equation}\label{cen:W-bound}
 \sqrt W\le\sqrt{V_0}+4.32\sqrt{qH_0}.
\end{equation}

The constants leave a strict margin for absorbing this spread:
\begin{align}
 \frac{p_*}{(1+p_*)qC_0}
 +\frac{p_*(4.32)^2}{1-p_*}
 &<\frac{.0532}{(.985)(.034)}
       +\frac{.019(4.32)^2}{.981}\label{cen:absorption-constant}\\
 &<1.95<2.\notag
\end{align}
For clarity, the middle rational number is
$444890776/228150625$, which is strictly less than $39/20$.
Weighted Cauchy--Schwarz, applied to \eqref{cen:W-bound}, now gives
\begin{align}
 \frac{p_*W}{2q}
 &\le\frac12\left(
       \frac{p_*}{(1+p_*)qC_0}
       +\frac{p_*(4.32)^2}{1-p_*}\right)
       \bigl((1+p_*)C_0V_0+(1-p_*)H_0\bigr)\notag\\
 &\le(1+p_*)C_0V_0+(1-p_*)H_0.
 \label{cen:absorption}
\end{align}
All denominators are positive by \eqref{cen:resources} and
$0\le p_*<.019$.

If \eqref{cen:variance-gain} is strict, substitute it into
\eqref{cen:tangent-energy}, then use \eqref{cen:spread} and
\eqref{cen:absorption}.  The result is
\[
 \mathcal E(d)>K+(1+p_*)C_0V_0+(1-p_*)H_0-p_*(q-l)\ge K,
\]
contradicting \eqref{cen:crossing}.  Lemma~\ref{cen:strict-variance}
therefore leaves only $m=0$ and $g^2=M$ everywhere.

In this exceptional case $b=s$, so $x_0=\sqrt{1-s^2}=\rho$.
Because $F$ is odd and strictly increasing, $g^2=M$ implies
$|d|=\rho$ everywhere.  If $W_k(f)$ denotes the Fourier weight of $f$
at degree $k$, Parseval and $m=0$ give
\[
 \sum_{k\ge1}W_k(f)=1,\qquad
 \sum_{k\ge1}r^kW_k(f)=\E d^2=r.
\]
Subtracting the second identity from $r$ times the first gives
$\sum_{k\ge2}(r-r^k)W_k(f)=0$.
Every coefficient $r-r^k$ is positive for $0<r<1$, so all weights above
degree one vanish.  Hence $f=\sum_i c_i x_i$, with $\sum_i c_i^2=1$.
The identity $f^2=1$ has degree-two coefficients $2c_ic_j=0$ for every
$i\ne j$.  At most one $c_i$ is nonzero, and its square is $1$.
Thus $f$ is a signed coordinate, proving Proposition~\ref{prop:central}.
\end{proof}
